% ========================================
% CAPÍTULO 2: FUNDAMENTACIÓN Y POSICIONAMIENTO  (criterio 2)
% ========================================

\chapter{FUNDAMENTACIÓN Y POSICIONAMIENTO}
\label{ch:fundamentacion}

La literatura se organiza en tres enfoques, que se comparan en la Tabla~\ref{tab:enfoques}.

\textbf{Detectores estáticos.} \textcite{yang2022} analizan el flujo de datos de notebooks para detectar patrones de fuga; \textcite{drobnjakovic2024} lo formalizan con interpretación abstracta; \textcite{alomar2025} lo llevan a una extensión de IDE. Son deterministas y reproducibles, pero ciegos a lo que depende de los valores de los datos, y se validaron sobre código humano.

\textbf{LLM como detector o juez.} \textcite{draoui2026} combinan \emph{slicing} estático con \emph{prompts} estructurados a un LLM, y \textcite{nam2025} incorporan un agente de chequeo de fuga dentro de un agente de ingeniería de ML. El LLM como juez tiene sesgos documentados \parencite{zheng2023} y no es reproducible, lo que lo hace inadecuado como instrumento de medición.

\textbf{Buenas prácticas.} \textcite{lones2024}, \textcite{kapoor2023} y la lista de verificación REFORMS \parencite{kapoor2024} catalogan la fuga y recomiendan encapsular el preprocesamiento; ese es el criterio de «código correcto» que usa este proyecto.

\begin{table}[H]
\centering
\caption{Enfoques existentes y posicionamiento del proyecto.}
\label{tab:enfoques}
\begingroup\footnotesize\singlespacing
\begin{tabular}{@{}p{2.9cm} p{2.7cm} p{2.6cm} p{2.2cm} p{4.6cm}@{}}
\toprule
Enfoque & Detecta & Mide el efecto sobre la métrica & Repara el código & Limitación principal \\
\midrule
Detectores estáticos & Patrones de flujo de datos & No & No & Ciegos a fuga dependiente de valores; validados en código humano \\
LLM como detector o juez & Lo descrito en el \emph{prompt} & No & A veces & Sesgos y variabilidad; sin evidencia trazable \\
\textbf{Este proyecto} & AST y conteo de filas en \texttt{.fit} & \textbf{Sí} (\DOpt, mismo pliegue) & \textbf{Sí} (patrones CV y Split) & Solo scikit-learn y tabular; alcance limitado de la reparación \\
\bottomrule
\end{tabular}
\endgroup
\end{table}

\textbf{Brecha y decisiones de diseño.} Entre las fuentes revisadas, ninguna mide el optimismo con una intervención pareada ni combina lectura y ejecución del código con reparación. De allí se derivan cuatro decisiones: (i) el \textbf{detector es determinista}, no un LLM, porque se usa como instrumento de medición y debe ser reproducible \parencite{zheng2023}; (ii) se combinan \textbf{AST y verificación en ejecución}, porque el primero lee el código pero no sus valores y el segundo observa lo que ocurrió; (iii) \textbf{\DOpt{} pareado} sigue la práctica de comparar una evaluación defectuosa con la corregida \parencite{kapoor2023,cawley2010}; (iv) la comparación entre familias usa pruebas no paramétricas pareadas \parencite{demsar2006}, porque con pocos datasets no se justifica asumir normalidad.
